% =============================================================================
% 02-os-operations-protection.tex — Lecture 2:
% Operating-System Operations and System Protection
%
% Course: Operating Systems (PCC CS-502)
% Anchor reading: Silberschatz, Galvin & Gagne, Operating System Concepts,
%                 8th ed. — Ch. 1 §1.5 (dual-mode operation, timer,
%                 protection, pp. 20-23); page-verified via
%                 master_supporting_docs/CS502/supporting_books/SilberschatzOSC/index.md.
%                 Stallings, Operating Systems: Internals and Design
%                 Principles, 6th ed. — Ch. 3 §3.5 (execution of the OS,
%                 p. 140) and §3.6 (security issues, p. 143); page-verified
%                 via .../StallingsOS/index.md.
%
% Compile with: /compile-latex CS502/02-os-operations-protection
% =============================================================================

\documentclass[aspectratio=169]{beamer}
\input{header}   % resolves via TEXINPUTS=../../Preambles;$TEXINPUTS — compile
                  % from cwd Slides/CS502/ (TeX writes output next to cwd)
\coursecode{CS502}

\title{Operating-System Operations and System Protection}
\subtitle{Lecture 2 --- Dual-Mode Operation, Privileged Instructions, and the Timer}
\author[Dr. Auqib Hamid Lone]{%
  \textbf{Dr. Auqib Hamid Lone}\\
  \small Assistant Professor in Computer Science \& Engineering
}
\institute[GCET Ganderbal]{%
  \small\textbf{PCC CS-502: Operating Systems}\\[0.3em]
  Department of Computer Science \& Engineering\\
  Government College of Engineering \& Technology (GCET), Ganderbal
}
\date[\today]{\small\today}

\begin{document}

\frame{\titlepage}

% =============================================================================
% ACT 0 --- THE CONTROL PROBLEM
% =============================================================================
\section{The Control Problem}

\begin{frame}{Welcome: The Program That Never Yields}
  \begin{block}{The big question}
    A student runs \texttt{while(1);} on a shared lab machine. The program
    executes no I/O, makes no system call, and never exits. Why does the
    machine not freeze --- and why can the instructor still kill the process?
  \end{block}
  \begin{itemize}
    \item The CPU has one core. Two programs cannot both be ``running'' at the
      same instant.
    \item If user code can execute \emph{every} instruction, nothing stops it
      from disabling interrupts, reprogramming the timer, or overwriting the
      kernel.
  \end{itemize}
  \vspace{0.15cm}
  \begin{alertblock}{Running thread for today}
    We follow one runaway process from its first instruction to the timer
    interrupt that takes the CPU away from it.
  \end{alertblock}
\end{frame}

\begin{frame}{Roadmap: Four Questions About Control}
  \begin{enumerate}
    \item \key{How does the CPU know who is running?} --- the mode bit.
    \item \key{What may user code not do?} --- privileged instructions.
    \item \key{How does the OS get the CPU back?} --- the timer and preemption.
    \item \key{How does a program ask for help?} --- the system-call gate.
  \end{enumerate}
  \vspace{0.2cm}
  \muted{Every answer is a hardware mechanism, not a software convention:
  protection only works because the hardware enforces it.}
\end{frame}

\transitionslide{The operating system is not a remote controller. It is a program that must run on the same CPU as the code it supervises.}

\begin{frame}{The OS Is a Program Too}
  \begin{itemize}
    \item The kernel uses the same CPU, the same registers, and the same
      physical memory as every user program.
    \item It cannot run ``in the background'' while user code runs; the CPU
      executes one instruction stream at a time.
    \item The kernel therefore runs only when an \key{event} transfers control
      to it: a system call, an interrupt, or an exception.
    \item The hardware must be able to tell, at every instant, whether the
      current instruction stream is \emph{trusted} kernel code or
      \emph{untrusted} user code.
  \end{itemize}
  \vspace{0.2cm}
  \begin{block}{The design problem}
    Build a machine that can run untrusted code at full speed, yet guarantee
    that the trusted kernel can always regain control.
  \end{block}
\end{frame}

\begin{frame}{What Goes Wrong Without a Boundary}
  \begin{itemize}
    \item A buggy pointer writes over another process's memory --- or over the
      kernel itself.
    \item A malicious program disables interrupts, and the timer can no longer
      fire.
    \item A program reprograms the timer to never expire, then loops forever.
    \item A program issues raw disk commands, bypassing the file system and
      every access-control check.
  \end{itemize}
  \vspace{0.2cm}
  \begin{alertblock}{The lesson}
    \bad{No boundary means no isolation, no multiplexing, and no trust.}
    The rest of this lecture is the mechanism that creates the boundary.
  \end{alertblock}
\end{frame}

% =============================================================================
% ACT 1 --- DUAL-MODE OPERATION
% =============================================================================
\section{Dual-Mode Operation}

\transitionslide{A single hardware bit separates the program that can do anything from the programs that must ask permission.}

\begin{frame}{One Bit, Two Worlds}
  \begin{block}{Dual-mode operation}
    A hardware \key{mode bit} records the current privilege level. The kernel
    runs in \key{kernel mode} (mode bit $=0$); user programs run in
    \key{user mode} (mode bit $=1$). Some instructions are legal only when the
    bit is $0$.
  \end{block}
  \begin{itemize}
    \item \textbf{Kernel mode:} full instruction set, access to all memory and
      all devices, can change the mode bit.
    \item \textbf{User mode:} restricted instruction set, access only to its own
      address space, cannot change the mode bit directly.
    \item The bit is \emph{hardware state}, not an ordinary memory variable.
  \end{itemize}
  \vspace{0.15cm}
  \muted{Silberschatz et al., Ch.~1 \S1.5, pp.~20--23.}
\end{frame}

\begin{frame}{Who May Execute What?}
  \begin{columns}[T]
    \begin{column}{0.48\textwidth}
      \begin{block}{Privileged (kernel only)}
        \begin{itemize}
          \item set or clear the mode bit
          \item disable or enable interrupts
          \item load the timer
          \item issue raw I/O instructions
          \item change the page-table base
          \item halt the machine
        \end{itemize}
      \end{block}
    \end{column}
    \begin{column}{0.48\textwidth}
      \begin{block}{Unprivileged (user or kernel)}
        \begin{itemize}
          \item arithmetic and logic
          \item loads and stores to the process's own memory
          \item branches and function calls
          \item read the clock (not set it)
          \item execute a trap instruction
        \end{itemize}
      \end{block}
    \end{column}
  \end{columns}
  \vspace{0.2cm}
  \muted{A privileged instruction executed in user mode does not silently
  fail: it \bad{traps} to the kernel, which normally terminates the process.}
\end{frame}

\begin{frame}{The Mode Transition Graph}
  \centering
  \begin{tikzpicture}[every node/.style={font=\small}]
    \node[dag-node, minimum width=3.4cm, minimum height=1.1cm] (user) at (-3.4,0)
      {User mode\\[2pt]\footnotesize mode bit $=1$};
    \node[dag-node, minimum width=3.4cm, minimum height=1.1cm] (kernel) at (3.4,0)
      {Kernel mode\\[2pt]\footnotesize mode bit $=0$};
    \draw[observed-edge] (user.east) --
      node[above, text=primary-blue] {return-from-trap} (kernel.west);
    \draw[observed-edge] (user.north) to[bend left=28]
      node[above, text=primary-blue] {trap / system call} (kernel.north);
    \draw[observed-edge] (user.south) to[bend right=28]
      node[below, text=primary-gold] {interrupt / exception} (kernel.south);
  \end{tikzpicture}
  \vspace{0.3cm}
  \begin{itemize}
    \item Entry into the kernel is always a \key{controlled} hardware event.
    \item The return path restores the saved user context and the user mode bit.
    \item The kernel does not ``call'' user code; it returns to it.
  \end{itemize}
\end{frame}

\begin{frame}{The Mode Bit Is Hardware, Not a Variable}
  \begin{itemize}
    \item If the mode bit were stored in ordinary memory, a user program could
      simply write \texttt{0} into it and become the kernel.
    \item The instruction that changes the bit is itself privileged, so user
      code cannot execute it.
    \item The hardware defines exactly which events may switch modes:
      interrupts, exceptions, and traps.
    \item This small, hardware-enforced boundary is the foundation of the
      \key{trusted computing base}.
  \end{itemize}
  \vspace{0.2cm}
  \begin{block}{Consequence}
    Protection is not a policy the OS can choose to enforce; it is a property
    the hardware guarantees. The OS decides \emph{what} to protect, but the
    mode bit decides \emph{whether} it can.
  \end{block}
\end{frame}

\begin{frame}{Worked Trace: One \texttt{write} Crosses the Boundary}
  \begin{enumerate}
    \item User code calls the C library \texttt{write(1, buf, n)}.
    \item The library places the system-call number and arguments in registers.
    \item It executes a \key{trap} instruction (\texttt{syscall} on x86-64).
    \item Hardware switches to kernel mode, saves the user PC and flags, and
      jumps to the trap handler.
    \item The kernel validates the arguments, performs the I/O, and stores a
      return value.
    \item A \key{return-from-trap} restores the user context and the user mode
      bit.
    \item The library returns the result; on failure it sets \texttt{errno}.
  \end{enumerate}
  \vspace{0.15cm}
  \muted{We will open this path up in Lecture 3. For now, notice that the mode
  switch is a hardware event, not a function call.}
\end{frame}

\begin{frame}{Socratic Check 1}
  \begin{block}{Question}
    Why can't the mode bit live in ordinary memory? Describe the attack that
    becomes possible if it does.
  \end{block}
  \vspace{0.25cm}
  \begin{itemize}
    \item A user program could store \texttt{0} into the mode variable.
    \item It could then disable interrupts, reprogram the timer, or overwrite
      the page tables.
    \item The hardware would no longer have a trusted state to return to.
  \end{itemize}
  \vspace{0.2cm}
  \muted{The bit must be outside the address space that user code can write.}
\end{frame}

% =============================================================================
% ACT 2 --- THE TIMER
% =============================================================================
\section{The Timer}

\transitionslide{Dual-mode protects memory and devices. It does not protect time: a user program can loop forever without executing a single privileged instruction.}

\begin{frame}{The Loop That Never Leaves User Mode}
  \begin{itemize}
    \item \texttt{while(1);} executes only unprivileged instructions.
    \item It never traps, never calls the OS, and never touches a device.
    \item If nothing external interrupts it, the kernel never runs again.
    \item The machine is not corrupted --- it is simply \bad{monopolized}.
  \end{itemize}
  \vspace{0.2cm}
  \begin{block}{The missing mechanism}
    The OS needs an event that fires \emph{without} the user program's
    cooperation. That event is the \key{timer interrupt}.
  \end{block}
\end{frame}

\begin{frame}{The Timer Is the External Event}
  \begin{itemize}
    \item A hardware timer counts down independently of the instruction stream.
    \item When it reaches zero, it raises a \key{timer interrupt}.
    \item The interrupt switches the CPU to kernel mode and jumps to the
      kernel's timer handler.
    \item The handler saves the running process's context and calls the
      scheduler.
    \item Loading the timer is a \key{privileged} instruction: user code cannot
      set it.
  \end{itemize}
  \vspace{0.2cm}
  \muted{The timer is a device, not a software alarm. If user code could reload
  it, preemption would be optional --- and a malicious program would opt out.}
\end{frame}

\begin{frame}{A Quantum on the Timeline}
  \centering
  \begin{tikzpicture}[x=1cm,y=1cm,every node/.style={font=\small}]
    \draw[->, thick, draw=jet] (0,0) -- (12.4,0) node[right] {time};
    \foreach \x/\lab in {0/0,3/10,6/20,9/30,12/40} {
      \draw[jet] (\x,0.08) -- (\x,-0.08);
      \node[below, text=jet] at (\x,0) {\lab\,ms};
    }
    \draw[fill=light-bg, draw=primary-blue, thick] (0,0.35) rectangle (3,0.85);
    \node at (1.5,0.6) {\texttt{P1}};
    \draw[fill=light-bg, draw=primary-gold, thick] (3,0.35) rectangle (6,0.85);
    \node at (4.5,0.6) {\texttt{P2}};
    \draw[fill=light-bg, draw=primary-blue, thick] (6,0.35) rectangle (9,0.85);
    \node at (7.5,0.6) {\texttt{P1}};
    \draw[fill=light-bg, draw=primary-gold, thick] (9,0.35) rectangle (12,0.85);
    \node at (10.5,0.6) {\texttt{P2}};
    \foreach \x in {3,6,9} {
      \draw[primary-gold, thick] (\x,0.9) -- (\x,1.25);
      \node[above, text=primary-gold] at (\x,1.25) {timer};
    }
  \end{tikzpicture}
  \vspace{0.35cm}
  \begin{itemize}
    \item Each process receives one 10\,ms \key{quantum}.
    \item The timer fires at the end of every quantum, even if the process is
      still runnable.
    \item The scheduler decides who runs next; the timer only guarantees that
      the decision happens.
  \end{itemize}
\end{frame}

\begin{frame}{Timer + Scheduler = Preemption}
  \begin{block}{Preemption}
    The kernel forcibly takes the CPU away from a running process at a
    timer interrupt, saves its context, and may schedule a different process.
  \end{block}
  \begin{itemize}
    \item \textbf{Cooperative:} a process yields voluntarily --- one infinite
      loop stops the system.
    \item \textbf{Preemptive:} the timer makes yielding unnecessary --- the
      kernel keeps control of the CPU.
  \end{itemize}
  \vspace{0.2cm}
  \muted{Trade-off: a large quantum wastes response time; a tiny quantum wastes
  CPU time on context switches.}
\end{frame}

\begin{frame}{Worked Example: What Does a Quantum Cost?}
  \begin{exampleblock}{Setup}
    A context switch costs $12\ \mu\text{s}$. The quantum is $Q$. What
    fraction of CPU time is spent switching?
  \end{exampleblock}
  \begin{center}
    \renewcommand{\arraystretch}{1.35}
    \begin{tabular}{c|c|c}
      \textbf{Quantum $Q$} & \textbf{Switches per second} & \textbf{Overhead} \\
      \hline
      $100\ \mu\text{s}$ & $10{,}000$ & $120{,}000\ \mu\text{s} = 12\%$ \\
      $1\ \text{ms}$     & $1{,}000$  & $12{,}000\ \mu\text{s} = 1.2\%$ \\
      $10\ \text{ms}$    & $100$      & $1{,}200\ \mu\text{s} = 0.12\%$ \\
    \end{tabular}
  \end{center}
  \vspace{0.1cm}
  \muted{Overhead $= \dfrac{12\ \mu\text{s}}{Q}$. The same formula explains why
  the quantum cannot simply be made as small as possible.}
\end{frame}

\begin{frame}{Socratic Check 2}
  \begin{block}{Question}
    If a small quantum improves response time, why not set the quantum to
    $1\ \mu\text{s}$?
  \end{block}
  \vspace{0.25cm}
  \begin{itemize}
    \item At $1\ \mu\text{s}$, the switch overhead alone is $1200\%$ --- the
      CPU spends more time switching than executing.
    \item Cache and TLB state are destroyed on every switch.
    \item The scheduler itself consumes CPU time.
    \item Real systems choose a quantum that balances responsiveness against
      overhead (often 1--100\,ms).
  \end{itemize}
\end{frame}

% =============================================================================
% ACT 3 --- PROTECTION IN DEPTH
% =============================================================================
\section{Protection in Depth}

\transitionslide{The mode bit and the timer protect the CPU. A complete boundary also protects memory and the devices.}

\begin{frame}{Three Things the Kernel Protects}
  \centering
  \begin{tikzpicture}[every node/.style={font=\small}]
    \node[dag-node, minimum width=3.2cm, minimum height=1.0cm] (cpu) at (-4.2,0)
      {CPU time};
    \node[dag-node, minimum width=3.2cm, minimum height=1.0cm] (mem) at (0,0)
      {Memory};
    \node[dag-node, minimum width=3.2cm, minimum height=1.0cm] (io) at (4.2,0)
      {I/O devices};
    \node[text=primary-gold] at (-4.2,-1.15) {timer};
    \node[text=primary-gold] at (0,-1.15) {MMU};
    \node[text=primary-gold] at (4.2,-1.15) {privileged I/O};
  \end{tikzpicture}
  \vspace{0.35cm}
  \begin{itemize}
    \item \textbf{CPU time:} the timer forces the kernel to run.
    \item \textbf{Memory:} the MMU checks every user access against the
      process's legal address range.
    \item \textbf{I/O:} devices are reachable only through privileged
      instructions and kernel drivers.
  \end{itemize}
\end{frame}

\begin{frame}{Memory Protection in One Slide}
  \begin{itemize}
    \item \textbf{Contiguous allocation:} base and limit registers bound the
      process; the hardware checks every address.
    \item \textbf{Paging:} the page table maps only the process's legal pages;
      an unmapped page traps to the kernel.
    \item \textbf{Segmentation:} base and limit registers per segment; the
      hardware checks the offset.
  \end{itemize}
  \vspace{0.2cm}
  \begin{block}{The common pattern}
    User code never sees a physical address. Every access is translated and
    checked by hardware; a violation becomes a \key{trap}, not a silent write.
  \end{block}
  \muted{Details in Unit III (Lectures 20--25).}
\end{frame}

\begin{frame}{I/O Protection}
  \begin{itemize}
    \item All raw I/O instructions are privileged: a user program cannot talk
      to a disk, network card, or timer directly.
    \item The user-level \texttt{read}/\texttt{write} calls are system calls;
      the kernel validates the buffer, the device, and the count.
    \item The kernel then invokes the device driver, which programs the
      controller.
    \item \texttt{cat} never touches the disk controller --- it asks the kernel,
      and the kernel decides whether the request is legal.
  \end{itemize}
  \vspace{0.2cm}
  \begin{alertblock}{Why this matters}
    Direct device access would let a process bypass file permissions, corrupt
    another process's data, or hang the device for everyone.
  \end{alertblock}
\end{frame}

\begin{frame}{Protection vs.\ Security}
  \centering
  \renewcommand{\arraystretch}{1.4}
  \begin{tabular}{p{0.42\textwidth}p{0.42\textwidth}}
    \textbf{Protection} & \textbf{Security} \\
    \hline
    Internal mechanisms among processes & External threats and attackers \\
    memory isolation, privilege levels, access control & authentication,
    malware defense, network firewalls \\
    Enforced by the OS and hardware & Enforced by policy, cryptography, and
    administrative practice \\
  \end{tabular}
  \vspace{0.35cm}
  \begin{itemize}
    \item Protection is a \emph{mechanism}; security is a \emph{policy} that
      uses the mechanism.
    \item A perfect mode bit does not help if the administrator grants every
      user root access.
  \end{itemize}
  \muted{Stallings, Ch.~3 \S3.6, p.~143.}
\end{frame}

\begin{frame}{The System-Call Gate}
  \centering
  \begin{tikzpicture}[every node/.style={font=\small}]
    \node[dag-node, minimum width=3.0cm] (user) at (-4.3,0) {User process};
    \node[dag-node, minimum width=2.6cm] (trap) at (-0.8,0) {Trap};
    \node[dag-node, minimum width=3.0cm] (table) at (3.0,0) {System-call table};
    \node[dag-node, minimum width=3.0cm] (svc) at (3.0,-2.0) {Kernel service};
    \draw[observed-edge] (user) -- (trap);
    \draw[observed-edge] (trap) -- (table);
    \draw[observed-edge] (table) -- (svc);
    \draw[observed-edge] (svc) -- ++(-3.6,0) -| (user);
  \end{tikzpicture}
  \vspace{0.25cm}
  \begin{itemize}
    \item The trap is the \key{only legal door} from user mode into the kernel.
    \item The kernel validates every argument and every pointer before acting.
    \item The return path restores user mode; the kernel never returns to
      arbitrary user code.
  \end{itemize}
  \muted{Lecture 3 opens this door and walks through it.}
\end{frame}

\begin{frame}{Socratic Check 3}
  \begin{block}{Question}
    If all I/O instructions are privileged, how does \texttt{cat} read a file?
    Trace the legal path in three steps.
  \end{block}
  \vspace{0.25cm}
  \begin{enumerate}
    \item \texttt{cat} calls the library \texttt{read}, which traps into the
      kernel.
    \item The kernel checks the file descriptor, buffer, and count, then asks
      the device driver to perform the transfer.
    \item The driver programs the controller; the kernel returns the data and
      restores user mode.
  \end{enumerate}
\end{frame}

% =============================================================================
% ACT 4 --- PUTTING IT TOGETHER
% =============================================================================
\section{Putting It Together}

\transitionslide{Protection is not one mechanism. It is a loop: the timer fires, the kernel runs, the scheduler chooses, and user code resumes --- or traps when it asks for something it cannot do.}

\begin{frame}{The Full Control Loop}
  \centering
  \begin{tikzpicture}[every node/.style={font=\small}]
    \node[dag-node, minimum width=3.0cm] (user) at (-4.0,-1.0) {User process running};
    \node[dag-node, minimum width=3.0cm] (kernel) at (-4.0,2.0) {Kernel: save context};
    \node[dag-node, minimum width=3.0cm] (sched) at (3.6,2.0) {Scheduler};
    \node[dag-node, minimum width=3.0cm] (restore) at (3.6,-1.0) {Kernel: restore context};
    \draw[observed-edge] (user.north) -- node[left, text=primary-gold] {timer} (kernel.south);
    \draw[observed-edge] (kernel.east) -- (sched.west);
    \draw[observed-edge] (sched.south) -- (restore.north);
    \draw[observed-edge] (restore.west) -- node[below, text=primary-gold] {return-from-trap} (user.east);
  \end{tikzpicture}
  \vspace{0.35cm}
  \begin{itemize}
    \item The loop runs on every timer tick and every blocking system call.
    \item The kernel is \emph{event-driven}: it runs because something happened,
      not because it is ``always on''.
    \item The user process cannot see the loop; it only sees time passing.
  \end{itemize}
\end{frame}

\begin{frame}{Worked Trace: Stopping the Runaway Process}
  \begin{enumerate}
    \item $t=0$: \texttt{P1} enters \texttt{while(1);} in user mode.
    \item $t=10$\,ms: the timer fires; hardware switches to kernel mode.
    \item The kernel saves \texttt{P1}'s PC and registers in its PCB.
    \item The scheduler picks \texttt{P2}; the kernel restores \texttt{P2}'s
      context and returns to user mode.
    \item $t=20$\,ms: the timer fires again; the scheduler may pick \texttt{P1}
      or a third process.
    \item If \texttt{P1} must be stopped, the kernel can deliver a signal or
      terminate it --- it already has control.
  \end{enumerate}
  \vspace{0.15cm}
  \muted{No cooperation from \texttt{P1} is required at any step.}
\end{frame}

\begin{frame}{Common Misconceptions}
  \begin{columns}[T]
    \begin{column}{0.48\textwidth}
      \begin{block}{Wrong}
        \begin{itemize}
          \item ``Dual-mode is a software convention.''
          \item ``The timer is a software alarm.''
          \item ``User code can disable interrupts if it is fast enough.''
          \item ``The OS runs continuously in the background.''
        \end{itemize}
      \end{block}
    \end{column}
    \begin{column}{0.48\textwidth}
      \begin{block}{Right}
        \begin{itemize}
          \item The mode bit and privileged instructions are hardware.
          \item The timer is a hardware device loaded by the kernel.
          \item Interrupt control is privileged; attempts trap.
          \item The kernel is event-driven, entered through traps and interrupts.
        \end{itemize}
      \end{block}
    \end{column}
  \end{columns}
\end{frame}

\begin{frame}{Takeaways}
  \begin{itemize}
    \item \key{Dual-mode operation} uses a hardware mode bit to separate
      trusted kernel code from untrusted user code.
    \item \key{Privileged instructions} protect the mode bit, the timer, the
      page tables, and the devices.
    \item \key{The timer} is the external event that lets the kernel preempt a
      runaway process.
    \item \key{The system-call gate} is the controlled entry point; the kernel
      validates every request.
    \item Protection is a \emph{hardware-enforced} boundary, not a promise the
      OS makes.
  \end{itemize}
\end{frame}

\begin{frame}{Next: System Calls}
  \begin{itemize}
    \item Today: how the OS stays in control --- mode bit, privileged
      instructions, timer, and the trap gate.
    \item Next lecture: what happens \emph{inside} the gate --- how a program
      names a service, how parameters travel, and what kinds of calls exist.
  \end{itemize}
  \vspace{0.3cm}
  \begin{block}{Reading}
    Silberschatz et al., \emph{Operating System Concepts}, Ch.~2
    \S2.3--2.4 \cite{SilberschatzOSC_operating_system_concepts}; Bovet \&
    Cesati, \emph{Understanding the Linux Kernel}, Ch.~10
    \cite{BovetCesati_understanding_linux_kernel}.
  \end{block}
\end{frame}

\begin{frame}{Further Reading (Page-Cited)}
  \footnotesize
  \begin{itemize}
    \item A.~Silberschatz, P.~B.~Galvin, G.~Gagne, \emph{Operating System
      Concepts}, 8e --- Ch.~1 \S1.5 (dual-mode operation, timer, protection,
      pp.~20--23).
    \item W.~Stallings, \emph{Operating Systems: Internals and Design
      Principles}, 6e --- Ch.~3 \S3.5 (execution of the OS, p.~140), \S3.6
      (security issues, p.~143).
    \item D.~P.~Bovet \& M.~Cesati, \emph{Understanding the Linux Kernel}, 3e
      --- Ch.~4 (interrupts and exceptions, pp.~131--188), Ch.~10 (system
      calls, pp.~398--419), for the concrete x86/Linux view.
  \end{itemize}
  \vspace{0.2cm}
  \muted{Page-verified via each book's \texttt{index.md} under
  \texttt{master\_supporting\_docs/CS502/supporting\_books/}.}
\end{frame}

\begin{frame}[allowframebreaks]{References}
  \footnotesize
  \bibliographystyle{plain}
  \bibliography{../../Bibliography_base}
\end{frame}

\end{document}
